% ============================================================
%  capitulo3_normativo.tex
%  Capítulo 3: Marco Normativo y Casos de Estudio en Colombia
% ============================================================

\chapter{Marco Normativo y Casos de Estudio en Colombia}
\label{chap:normativo}

El marco normativo que regula la normalización de redes eléctricas en Colombia lleva más de tres décadas construyéndose, respondiendo a aprendizajes de campo, a las necesidades del sector y, más recientemente, a los compromisos del PND 2022--2026 con la transición energética. Entender esa trayectoria ayuda a explicar por qué el esquema actual tiene las oportunidades que tiene, pero también los vacíos que tiene.

Este capítulo cubre tres bloques: la regulación asociada a la normalización de redes, la regulación de los sistemas AGPE, y los casos de estudio documentados en Colombia. Los tres son necesarios para el análisis técnico del Capítulo~\ref{chap:tecnico} y para la metodología del Capítulo~\ref{chap:metodologia}.


% ─────────────────────────────────────────────────────────────
\section{Evolución Normativa Asociada a la Normalización de Redes Eléctricas}
\label{sec:normativa_normalizacion}
% ─────────────────────────────────────────────────────────────

\subsection{Las leyes fundacionales del sector eléctrico colombiano (1994)}
\label{subsec:leyes_1994}

El punto de partida es 1994. Las \textbf{Leyes 142 y 143} establecieron los principios sobre los que todavía opera el sector eléctrico colombiano \parencite{Ley142_1994, Ley143_1994}.

La \textbf{Ley 142} definió los criterios de prestación de los servicios públicos domiciliarios: eficiencia, continuidad, calidad, solidaridad. En lo que respecta a los BS, su artículo 9 es claro: los usuarios tienen derecho a recibir un servicio medido con instrumentos adecuados y a que el consumo facturado sea real. El artículo 11 obliga a las empresas prestadoras a facilitar a los usuarios de menores ingresos el acceso a los subsidios disponibles. Esos subsidios cubren el 50\%, 40\% y 15\% del consumo básico de subsistencia para los estratos 1, 2 y 3 respectivamente, financiados con aportes del 20\% de los estratos 5, 6 y el sector comercial e industrial \parencite{Ley142_1994}.

La \textbf{Ley 143} complementó este marco para el sector eléctrico específicamente. Su artículo 11 definió el consumo básico de subsistencia como la cantidad mínima de electricidad mensual que un usuario típico necesita para satisfacer sus necesidades básicas. El artículo 48 ordenó al Gobierno Nacional destinar recursos del PND a programas de energización prioritarios para lograr cobertura igualitaria en el país. El artículo 66 estableció el ahorro y la eficiencia energética como objetivos del sector \parencite{Ley143_1994}.

\subsection{La creación del PRONE (2003--2006)}
\label{subsec:creacion_prone}

El PRONE nació en el PND 2003--2006, aprobado por la \textbf{Ley 812 de 2003}, que en su artículo 63 lo creó como instrumento de financiamiento para la normalización de redes eléctricas \parencite{Ley812_2003}. Tres años después, la \textbf{Ley 1117 de 2006} le dio soporte legal definitivo y lo constituyó formalmente como un fondo administrado por el MME, con el propósito de financiar planes, programas y proyectos para legalizar usuarios\footnote{La legalización es el proceso de individualizar la medida, crear el código de usuario en el sistema del OR y establecer la relación contractual con factura individual.} y adecuar las redes en los BS del SIN \parencite{Ley1117_2006}.

La reglamentación detallada llegó con el \textbf{DUR 1073 de 2015}, artículo 2.2.3.3.3.1.1, que delimitó los rubros financiables dentro del programa: \textit{``únicamente el suministro e instalación de las redes de distribución, los transformadores de distribución, las acometidas a las viviendas de los usuarios y los medidores o sistema de medición del consumo. En lo referente al desmonte del material existente [\ldots] su costo no podrá superar el tres por ciento (3\%) del valor total del proyecto''} \parencite{DUR1073_2015}.

El DUR también define quiénes pueden presentar proyectos, los requisitos mínimos, los criterios de priorización y las condiciones de propiedad de los activos una vez terminadas las obras. Es el marco institucional vigente dentro del cual se rigen los OR que buscan financiamiento público para la normalización de redes eléctricas.

Adicionalmente, consolida la reglamentación del \textbf{FOES}, subsidio para usuarios de estratos 1 y 2 en áreas especiales. Su aplicación depende del consumo de subsistencia definido por la UPME, del porcentaje de recaudo del OR —la cobertura del subsidio varía entre el 40\% y el 100\% según ese indicador— y de la certificación anual del área especial ante la SSPD \parencite{DUR1073_2015}.

El mecanismo de financiamiento del PRONE quedó fijado por la \textbf{Ley 1753 de 2015}: el fondo recibe \$1,90 por kilovatio hora transportado en el SIN a través del ASIC. El FOES, por su parte, cubre hasta \$92 por kilovatio hora destinado al consumo de subsistencia de los estratos 1 y 2 \parencite{Ley1753_2015}.

\subsection{Los esquemas diferenciales de facturación}
\label{subsec:esquemas_diferenciales}

El DUR 1073 reconoce algo que en la práctica resulta evidente: las reglas convencionales de facturación no funcionan en los BS. Por eso establece esquemas alternativos que los OR pueden usar: medición y facturación comunitaria, proyecciones de consumo basadas en cargas contratadas o históricos, pago anticipado y periodos de facturación flexibles \parencite{DUR1073_2015}.

Para aplicar cualquiera de estos esquemas, el OR debe firmar contratos con suscriptores comunitarios que definan cómo se mide, cómo se factura y cómo se paga. El problema práctico es que el suscriptor comunitario no tiene herramientas legales reales para exigir el pago a los usuarios individuales.

\subsection{La Resolución CREG 015 de 2018 y la remuneración de la distribución}
\label{subsec:resolucion015}

La \textbf{Resolución 015 de 2018} de la CREG estableció la metodología para remunerar la actividad de distribución en el SIN \parencite{ResolucionCREG015_2018}. Los capítulos 6 y 7 —sobre planes de inversión y gestión de pérdidas— son los más relevantes para este trabajo: definen cómo se remunera vía tarifa la inversión orientada a reducir pérdidas, mejorar calidad y aumentar demanda.

\subsection{La incorporación del AGPE al PRONE (2023--2024)}
\label{subsec:agpe_prone}

El cambio normativo más importante en los veinte años de historia del PRONE llegó con el \textbf{PND 2022--2026}, aprobado por la \textbf{Ley 2294 de 2023}: por primera vez, los recursos del programa pueden destinarse a financiar sistemas AGPE dentro de los proyectos de normalización \parencite{Ley2294_2023}.

La reglamentación vino con el \textbf{Decreto 1023 de 2024}, que adicionó la Subsección 3.4 al DUR 1073. El decreto fijó los requisitos para presentar proyectos con AGPE, los criterios de priorización y las reglas de propiedad de los activos \parencite{Decreto1023_2024}. Una novedad relevante: los activos de generación construidos con recursos del PRONE pueden transferirse a los usuarios beneficiados a través de la figura de \textbf{Comunidades Energéticas}, reglamentadas por el \textbf{Decreto 2236 de 2023}, que les otorga carácter legal de prestador de servicio público \parencite{Decreto2236_2023, DUR1073_2015}.

El \textbf{Decreto 1403 de 2024} amplió las posibilidades técnicas al permitir que un sistema AGPE entregue excedentes al SIN o consuma energía en sitios distintos a los de producción \parencite{Decreto1403_2024}.

La primera convocatoria con el AGPE como gasto elegible fue la \textbf{Resolución 40243 de 2025} \parencite{ResolucionMME40243_2025}. El resultado fue elocuente: ninguno de los proyectos presentados cumplió con todos los requisitos del MME \parencite{MME2025}. Ese dato confirma, de la manera más directa posible, la brecha metodológica que esta monografía busca cerrar.


% ─────────────────────────────────────────────────────────────
\section{Evolución Normativa Asociada a los Sistemas de AGPE}
\label{sec:normativa_agpe}
% ─────────────────────────────────────────────────────────────

\subsection{Los orígenes restrictivos de la autogeneración (1994--2014)}
\label{subsec:origenes_agpe}

La \textbf{Ley 143 de 1994} definió al autogenerador como \textit{``aquel generador que produce energía eléctrica exclusivamente para atender sus propias necesidades''} \parencite{Ley143_1994}. Era autoconsumo puro, sin posibilidad de excedentes, sin red comunitaria y sin articulación con la política energética.

La \textbf{Resolución 84 de 1996} reglamentó eso con un esquema aún más restrictivo: el autogenerador conectado al SIN no podía vender excedentes, salvo en situaciones de racionamiento y solo a través de la Bolsa de Energía. La conexión al sistema se concebía exclusivamente como respaldo. Para los usuarios regulados, el suministro iba obligatoriamente a través del comercializador del área; los no regulados podían contratar libremente.

\subsection{La Ley 1715 de 2014: el cambio estructural}
\label{subsec:ley1715}

La \textbf{Ley 1715 de 2014} cambió las reglas de fondo \parencite{Ley1715_2014}. Incorporó las FNCER a la política energética nacional y redefinió la autogeneración: dejó de ser exclusivamente autoconsumo y pasó a reconocerse como una alternativa con posibilidad de entregar excedentes a la red bajo condiciones que definiría la CREG.

Los beneficios que trajo son concretos: incentivos tributarios para proyectos con FNCER (exención de IVA, deducción de renta, exención de aranceles), reconocimiento de la autogeneración como instrumento de eficiencia energética y posibilidad de medición bidireccional. Ese último punto es técnicamente fundamental: sin medición bidireccional no hay forma de contabilizar lo que el sistema inyecta a la red.

Entre 2014 y 2018 la normativa reglamentaria fue tomando forma. La \textbf{Resolución UPME 281 de 2015} fijó el límite de 1~MW para los AGPE \parencite{ResolucionUPME281_2015}. El \textbf{Decreto 348 de 2017} simplificó las condiciones de medición, conexión y contrato de respaldo \parencite{Decreto348_2017}. La \textbf{Resolución CREG 030 de 2018} fue el primer intento integral de regular los AGPE en el SIN \parencite{ResolucionCREG030_2018}, pero generó tantas preguntas de aplicación que la CREG tuvo que revisarla.

\subsection{La Resolución CREG 174 de 2021: la norma vigente}
\label{subsec:resolucion174}

La \textbf{Resolución CREG 174 de 2021} derogó completamente la 030 de 2018 y es hoy la referencia regulatoria para los sistemas AGPE en el SIN \parencite{ResolucionCREG174_2021}. Tres cambios son particularmente relevantes para proyectos en BS:

\begin{itemize}
    \item \textbf{El límite de inyección subió del 15\% al 50\%} de la capacidad del transformador o circuito al que se conecta el AGPE. Para los circuitos que sirven a los BS, que generalmente tienen capacidades nominales relativamente bajas, el umbral anterior del 15\% limitaba significativamente la potencia instalable del sistema AGPE sin requerir refuerzos en la red. Con el nuevo límite del 50\%, la potencia de generación que puede conectarse al mismo punto aumenta considerablemente, lo que mejora la viabilidad técnica y económica de estos proyectos al permitir sistemas de mayor tamaño sin inversiones adicionales en infraestructura de red.

    \item \textbf{La información de conexión se hizo más accesible.} La resolución obligó a los OR a mantener un micrositio específico en su página web donde cualquier interesado puede consultar las características técnicas del punto de conexión y el estado de la red, sin software adicional. Eso reduce una barrera real para quien está formulando un proyecto.

    \item \textbf{Se definió el manejo de excedentes vía tarifa}, distinguiendo el caso en que los excedentes inyectados a la red superan la energía importada desde ella (denominado permuta: el AGPE entrega más energía de la que toma, y el OR compensa al usuario a una tarifa definida por la CREG) del caso en que los excedentes son menores o iguales a las importaciones (el excedente se descuenta directamente de la factura). Esa distinción alinea la regulación con la realidad técnica de la medición bidireccional y define cómo se liquida el intercambio de energía entre el usuario y el OR.
\end{itemize}

\subsection{La integración del AGPE en la normalización (2023--2025)}
\label{subsec:integracion_agpe_normalizacion}

El Decreto 1023 de 2024 es la pieza normativa que hace posible lo que propone esta monografía. Integra el AGPE dentro del PRONE y crea la distinción de propiedad que ya se mencionó: la red normalizada queda bajo el OR para efectos de AOM; los activos de generación pueden transferirse a la Comunidad Energética. Esa distinción tiene consecuencias sobre quién responde por qué en el largo plazo, y requiere que el proyecto la defina con claridad desde el inicio \parencite{Decreto1023_2024}.

El costo unitario por usuario en proyectos con AGPE es notablemente superior al de proyectos convencionales. Los 30 proyectos de la vigencia 2025 lo ilustran: requirieron un valor de financiación significativamente mayor al de vigencias anteriores con muchos más proyectos \parencite{MME2025}. Esa presión sobre los recursos disponibles del PRONE refuerza la necesidad de criterios de priorización más específicos para proyectos con AGPE, uno de los vacíos regulatorios que se aborda en el Capítulo~\ref{chap:metodologia}.


% ─────────────────────────────────────────────────────────────
\section{Análisis de la Normalización Eléctrica en Colombia: Experiencias, Financiación y Cuellos de Botella}
\label{sec:casos_estudio}
% ─────────────────────────────────────────────────────────────

Tres décadas de proyectos de normalización en Colombia dejan un registro valioso. Esta sección analiza los tres casos con mayor documentación disponible en la literatura, y, de manera complementaria, se analizó el comportamiento histórico de los proyectos del PRONE entre 2005 y 2025.

\subsection{Caso 1 -- Cartagena, Bolívar (1991)}
\label{subsec:caso_cartagena}

En 1991, la Electrificadora de Bolívar evaluó la viabilidad de construir redes de distribución en zonas subnormales de Cartagena usando el método de normalización por tramos secundarios, siendo uno de los primeros estudios documentados de este tipo en el país \parencite{CastilloQuintana1991}.

El diseño propuesto era sencillo: un vano de 80~m de red primaria (13,2~kV) y dos vanos de 40~m de red secundaria (220~V), con postes de 12~m y 8~m ubicados en calles principales. Se distinguía entre usuarios directos —más cerca de la calle, con alumbrado público, facturación sobre el 100\% del consumo— e indirectos —más alejados, sin alumbrado, con recaudo sobre el 50\% a través de la JAC—. Se esperaba beneficiar a 24 usuarios directos y 18 indirectos. La financiación proyectada era mediante créditos del BID y la FEN, con recuperación de inversión en 18 meses.

El caso tiene varias limitaciones estructurales que, tres décadas después, siguen siendo completamente relevantes. El modelo colectivo dividía el consumo por partes iguales sin reconocer el ahorro individual, lo que no incentivaba el uso racional de la energía. La brecha entre usuarios directos e indirectos generaba inequidades concretas en acceso al alumbrado y disponibilidad de potencia. Y los usuarios ya entonces declaraban destinar el 80\% de sus ingresos a necesidades básicas, lo que ponía en duda la capacidad real de pago sobre la que se sostenía el modelo financiero \parencite{CastilloQuintana1991}.

\subsection{Caso 2 -- Sabana de Torres, Santander (2010--2014)}
\label{subsec:caso_santander}

La Electrificadora de Santander (ESSA) ejecutó entre 2010 y 2014 uno de los proyectos de normalización más completos ejecutados en el país, orientado a reducir pérdidas no técnicas que en 2009 representaban aproximadamente el 20\% de la energía transportada \parencite{MendezFuentes2013}.

La inversión fue íntegramente con recursos propios del OR, distintos de los recursos públicos del Ministerio que provienen del presupuesto general de la nación: \$114.885 millones en total, de los cuales \$38.894 millones fueron inversión directa en infraestructura y \$75.991 millones costos operativos. La estrategia técnica combinó tres componentes:

\begin{enumerate}[label=\roman*), leftmargin=2cm, itemsep=4pt]
    \item \textbf{Red Chilena:} En lugar de red de baja tensión convencional, se instalaron transformadores pequeños de 5 a 15~kVA que alimentan directamente entre 6 y 9 usuarios a través de la acometida, desde una caja de derivación. Sin red de BT, no hay dónde conectarse ilegalmente.

    \item \textbf{Conductor blindado antifraude:} Se sustituyó el trenzado por cable concéntrico TRI y TETRA CAPA, que no puede ser perforado con los conectores informales que se usaban para las conexiones clandestinas. El conductor del alumbrado público también se redujo de calibre (a cable 2×18) para que solo tuviera capacidad para las luminarias.

    \item \textbf{AMI en modalidad prepago:} Los medidores no se instalaron dentro de las viviendas para evitar manipulaciones; solo se dejaron displays de visualización en el interior. El objetivo era desincentivar el fraude sin perder la transparencia hacia el usuario.
\end{enumerate}

Al cierre de 2013 se habían ejecutado 38.285 legalizaciones y normalizado 110.064 usuarios en total \parencite{MendezFuentes2013}. Los resultados son sólidos. Sin embargo, la lección más importante es sobre escala: una inversión de \$114.885 millones con recursos propios está fuera del alcance de la mayoría de los OR del país, especialmente de los que operan en el Caribe con los problemas financieros documentados en el Capítulo~\ref{chap:diagnostico}. Vale precisar que el costo por usuario del PRONE actualmente se estima alrededor de \$11 millones, lo que representa un esquema de financiamiento más accesible para los OR de menor tamaño. Sin el PRONE, este tipo de proyectos difícilmente sería replicable en el contexto caribeño.

\subsection{Caso 3 -- Popayán, Cauca (2014--2018)}
\label{subsec:caso_popayan}

La experiencia de la Compañía Energética de Occidente (CEO) en Popayán se distingue por la integralidad del proceso: combinó intervención técnica, gestión social sistemática y articulación interinstitucional en un solo proyecto \parencite{SerranoRueda2019}.

En cuatro años se normalizaron 41 asentamientos en Popayán —establecidos por más de 15 años— con 12.000 habitantes. El financiamiento fue mixto: el OR aportó el 76\% del valor total (\$3.500 millones) y el MME el 24\% restante a través del PRONE. Esa estructura muestra algo importante: el PRONE no necesita ser el financiador principal para ser determinante; puede funcionar como el factor que hace viable lo que el OR solo no podría asumir.

La estrategia tuvo tres componentes. El técnico incluyó AMI y modernización completa de la infraestructura de distribución secundaria. El social fue el más intenso: formación comunitaria sistemática sobre derechos y deberes, lectura de factura y uso eficiente de la energía, realizada antes, durante y después de las obras. El componente de eficiencia energética fue el ``Cambiatón'': sustitución masiva de luminarias incandescentes por ahorradoras en las viviendas de los usuarios, para que la transición a la facturación individual no llegara acompañada de un consumo sin reducir.

Los resultados hablan solos: recaudo del 0\% al 97\%, reducción del 67\% en el consumo por usuario (de 197 a 65~kWh-mes), mejora en calidad del servicio y condonación de 15 años de deudas acumuladas. Lo que no se puede ignorar es que articular todo eso requirió más de diez entidades — alcaldía, personería, defensoría, concejo, procuraduría, SSPD, secretarías, fuerza pública y medios de comunicación. Sin esa articulación, los asentamientos con mayor resistencia social no se hubieran podido intervenir \parencite{SerranoRueda2019}.

\subsection{Comportamiento histórico del PRONE (2005--2025)}
\label{subsec:prone_historico}

A partir de las actas CAPRONE y los informes de evaluación del MME se construyó una base de datos de los proyectos adjudicados por vigencia, relacionándose los OR seleccionados, ubicación del proyecto, número de beneficiarios y valor solicitado \parencite{MME2025}. De esta se estableció que la concentración de proyectos financiados por el MME fueron presentados para la región Caribe, Pacífica y algunos departamentos de la andina. En cuanto a Orinoquía y Amazonía la participación fue marginal, lo cual es coherente con la existencia de BS documentados en esa zona.

El 52,03\% de los proyectos presentados entre 2005 y 2025 provienen de la Electrificadora de la Costa Caribe (hoy Afinia y Air-e); en 2024, ningún proyecto de esos OR —ni de ningún otro— resultó elegible.

Entre las variaciones anuales se destacan las siguientes:

\begin{itemize}
    \item \textbf{2007--2010}: Sin proyectos aprobados. El ajuste normativo tras la Ley 1117 de 2006 y el Decreto 1123 de 2008 redefinió actores, condiciones y procedimientos, y el programa estuvo prácticamente detenido.

    \item \textbf{2012}: 185 proyectos, el pico histórico del programa. Probablemente el resultado de años de acumulación de diagnósticos técnicos en los OR.

    \item \textbf{2016--2018}: De nuevo sin proyectos. El DUR 1073 de 2015 generó otro periodo de ajuste institucional.

    \item \textbf{2024}: Cero proyectos elegibles. El primer año con convocatoria bajo el nuevo esquema AGPE y ninguno cumplió los requisitos.

    \item \textbf{2025}: 30 proyectos con el esquema nuevo, beneficiando alrededor de 11.000 usuarios en La Guajira, Magdalena, Atlántico, Nariño, Huila y Meta. El valor de financiación solicitado fue significativamente mayor al de vigencias con mucho más proyectos, por el costo del componente solar.
\end{itemize}

En total, entre 2003 y 2025 el PRONE ha beneficiado aproximadamente a 290.000 usuarios según los proyectos aprobados. Esa cifra puede diferir de los beneficiarios reales por los ajustes que ocurren durante la ejecución.

\begin{table}[H]
\centering
\caption{Casos de estudio de normalización en Colombia con recursos públicos y privados}
\label{tab:casos_estudio}
\begin{tabularx}{\textwidth}{>{\RaggedRight}p{2.5cm}
                             >{\RaggedRight}p{2.8cm}
                             >{\RaggedRight}p{3.2cm}
                             >{\RaggedRight\arraybackslash}p{4.3cm}}
\toprule
\textbf{Caso de Estudio} & \textbf{Financiación / Normatividad} & \textbf{Impacto en los Usuarios} & \textbf{Impacto en el OR} \\
\midrule
\rowcolor{grisclaro}
\textbf{Cartagena, Bolívar} (1991)
& Créditos BID-FEN. Sin regulación sobre PR ni esquemas diferenciales.
& \textit{Negativo:} Usuarios asumían costos de medición y todas las pérdidas. Brecha entre usuarios directos e indirectos.
& \textit{Cuello de botella:} Recaudo delegado a JAC sin herramientas legales. \newline\textit{Positivo:} Proyección de recuperación a 18 meses. \\

\textbf{Sabana de Torres, Santander} (2010--2014)
& Recursos propios del OR (\$114.885 millones).
& \textit{Positivo:} Información en tiempo real (AMI), infraestructura blindada, eliminación de conexiones fraudulentas.
& \textit{Negativo:} Alta inversión, reincidencia en sectores sin seguimiento. \newline\textit{Positivo:} Reducción de pérdidas con Red Chilena. \\

\rowcolor{grisclaro}
\textbf{Popayán, Cauca} (2014--2018)
& Mixta: OR 76\% + PRONE 24\%. Articulación con MME.
& \textit{Positivo:} Consumo redujo 67\%, deudas históricas condonadas, mejora en seguridad y sentido de legalidad.
& \textit{Positivo:} Recaudo 0\% a 97\%, mejora en calidad. \newline\textit{Negativo:} Asumir cartera morosa. \newline\textit{Cuello de botella:} Articulación con más de 10 instituciones. \\
\bottomrule
\multicolumn{4}{l}{\footnotesize\textit{Fuente: Elaboración propia con datos de \parencite{CastilloQuintana1991, MendezFuentes2013, SerranoRueda2019}.}}
\end{tabularx}
\end{table}

\subsection{Lo que enseñan los tres casos}
\label{subsec:sintesis_casos}

Cuatro patrones se repiten con independencia del OR, el periodo o la región:

\textbf{La tecnología no se sostiene sola.} Santander y Popayán muestran que la infraestructura avanzada mejora los resultados técnicos, pero si la comunidad no la apropia, los usuarios encuentran la manera de sortearla. El blindaje más sofisticado no reemplaza al acompañamiento social.

\textbf{La gestión social previa es lo que diferencia el éxito del fracaso.} No el diseño eléctrico. El contraste entre Popayán —con formación comunitaria sistemática— y los proyectos cancelados en el Caribe por oposición organizada lo ilustra de manera directa.

\textbf{El recaudo no mejora solo con la individualización.} La normalización crea las condiciones para mejorar el recaudo, pero no lo garantiza. Depende de la capacidad de pago de los usuarios, de los mecanismos de cobro del OR y de si se condona o no la deuda histórica.

\textbf{Los cuellos de botella más frecuentes son institucionales, no técnicos.} La falta de certificación del BS, la debilidad de algunos OR para estructurar proyectos elegibles y la complejidad de los requisitos del PRONE no se resuelven con mejores equipos de medición. Requieren herramientas metodológicas claras, que es precisamente lo que propone el Capítulo~\ref{chap:metodologia}.